\chapter{Acreditación}
\section{¿Qué es la acreditación}
La acreditación es la herramienta establecida a escala internacional para \textbf{generar confianza} sobre la actuación de un tipo de organizaciones muy determinado que se denominan de manera general \textbf{Organismos de Evaluación de la Conformidad} y que abarca a los \textit{Laboratorios de ensayo, Laboratorios de Calibración, Entidades de Inspección, Entidades de certificación y Verificadores Ambientales}

\section{¿Por qué existe?}
Para lograr la \textbf{CONFIANZA Y CREDIBILIDAD} que el mercado y la sociedad en general tiene de los \textbf{evaluadores que realizan las evaluaciones}, por lo que es preciso establecer un mecanismo independiente, riguroso y global que \textbf{garantice la competencia técnica} de dichos evaluadores y su sujeción a normas de carácter internacional.

\section{Entidades de Acreditación}
\textit{``Se reconoce y designa a la Entidad Nacional de Acreditación (ENAC) como Entidad de Acreditación de las establecidas en el Capítulo 2''.} - RD 2200:1995 (modificado por el RD 338:2010)

\paragraph{CAPÍTULO 2. ENTIDADES DE ACREDITACIÓN.} \textit{``Son entidades privadas sin ánimo de lucro, que se constituyen con la finalidad de acreditar en el ámbito estatal, a las entidades de certificación, laboratorios de ensayo y calibración y entidades auditoras y de inspección que actúen en el ámbito voluntario de la CALIDAD, así como a los Organismos de Control que actúen en el ámbito reglamentario y a los Verificadores medioambientales, mediante la verificación del cumplimiento de las condiciones y requisitos técnicos exigidos para su funcionamiento.''}

\section{¿Por qué existen los Evaluadores de la Conformidad?}
El objetivo principal de la actuación de los organismos de evaluación de la conformidad es el de demostrar a la sociedad (Autoridades, empresas y consumidores en general) que los PRODUCTOS Y SEVRICIOS puestos a su disposición son conformes con ciertos requisitos relacionados generalmente con su CALIDAD y la SEGURIDAD.

Dichos requisitos pueden estar establecidos por la Ley y tneer por tanto CARÁCTER REGLAMENTARIO o estar especificados en Normas, especificaciones u otros documentos de CARÁCTER VOLUNTARIO.

\section{¿Cuál es el proceso de una acreditación?}

